\documentclass[
    left=2.5cm,         % Sadly, generic margin parameter
    right=2.5cm,        % doesnt't work, as it is
    top=2.5cm,          % superseded by more specific
    bottom=3cm,         % left...bottom parameters.
    bindingoffset=6mm,  % Optional binding offset.
    nohyphenation=false % You may turn off hyphenation, if don't like. 
]{eiti/eiti-thesis}

\usepackage[polish]{babel}
\usepackage[
    backend=bibtex,
    style=ieee
]{biblatex}
\usepackage{csquotes}

\graphicspath{{img/}}             % Katalog z obrazkami.
\addbibresource{bibliografia.bib} % Plik .bib z bibliografią

%----------------------------------------
% Twierdzenia i definicje;
% tutaj ew. tłumaczymy te terminy
% na inne języki
%----------------------------------------
\newtheorem{theorem}{Twierdzenie}
\newtheorem{lemma}{Lemat}
\newtheorem{corollary}{Wniosek}
\newtheorem{definition}{Definicja}
\newtheorem{axiom}{Aksjomat}
\newtheorem{assumption}{Założenie}

%----------------------------------------
% Spis rysunków, tablic i załączników;
% tutaj ew. tłumaczymy te terminy
% na inne języki
%----------------------------------------
\AtBeginDocument{
    \renewcommand{\listfigurename}{Spis rysunków}
    \renewcommand{\listtablename}{Spis tabel}
    \renewcommand{\tablename}{Tabela}
}

\begin{document}

%--------------------------------------
% Strona tytułowa
%--------------------------------------
\EngineerThesis % dla pracy inżynierskiej mamy \EngineerThesis
\instytut{Automatyki i Informatyki Stosowanej}
\kierunek{Automatyka i Robotyka}

\title{
     System rozproszonej i bezpiecznej identyfikacji użytkownika urządzenia Internetu Rzeczy
}
\engtitle{ % Tytuł po angielsku do angielskiego streszczenia
    System for distributed and secure identification for the user of IoT device
}

\author{Dawid Kiciński}
\album{277115}
\promotor{mgr Maciej Stefańczyk}
\date{\the\year}
\maketitle

%--------------------------------------
% Streszczenie po polsku
%--------------------------------------
\streszczenie \lipsum[1-3]
\slowakluczowe XXX, XXX, XXX
\newpage

%--------------------------------------
% Streszczenie po angielsku
%--------------------------------------
\abstract \kant[1-3]
\keywords XXX, XXX, XXX
\newpage

%--------------------------------------
% Oświadczenie o autorstwie
%--------------------------------------
\makeauthorship
\blankpage

%--------------------------------------
% Spis treści
%--------------------------------------
\thispagestyle{empty}
\tableofcontents
\blankpage

%--------------------------------------
% Rozdziały
%--------------------------------------
\newpage
\section[Wstęp]{Wstęp}
To jest jakiś tekst po polsku. 

\section{Rozproszone uczenie na urządzeniu końcowym}

\subsection{Algorytm aktualizacji modelu globalnego}


\section{System weryfikacji użytkownika}

Zadaniem systemu jest weryfikacja użytkownika, cz zostały przedstiawione w następnej sekcji.

\subsection{Wstępne przetwarzanie obrazu}

\subsection{Weryfikacja użytkownika}
Weryfikacja twarzy jest zadaniem przyrównania twarzy kandydata to innej, i weryfikacja czy nastąpiło ich dopasowanie. Jest to mapowanie jeden-do-jednego: należy sprawdzić czy jest to ta sama osoba.

\subsubsection{Procedura weryfikacji}
\subsubsection{Ekstraktor cech}


% \newpage % Zaleca się otwieranie rozdziału od nowej strony.
\section[Wstęp]{Wstęp}
\lipsum[1] \cite{goossens93}
\begin{figure}[!h]
    \label{fig:anzelm}
    \centering \includegraphics[width=0.5\linewidth]{logopw.png}
    \caption{Tradycyjne godło Politechniki Warszawskiej.}
\end{figure}
\lipsum[2-10]
         % W długich pracach
% \newpage % Zaleca się otwieranie rozdziału od nowej strony.
\section{De Finibus Bonorum et Malorum}
\lipsum[1] Lorem ipsum dolor sit amet\footnote{Lorem ipsum dolor sit amet, consectetur adipiscing elit, sed do eiusmod tempor incididunt ut labore et dolore magna aliqua. Ut enim ad minim veniam, quis nostrud exercitation ullamco laboris nisi ut aliquip ex ea commodo consequat.}. 
\begin{align*}
    E & = mc^2 \\ 
    y & = ax^2 + bx + c
\end{align*}

\lipsum[3]

\begin{align}
\begin{bmatrix}
    1 & 0 & 0 \\ 
    0 & 2 & 0 \\ 
    0 & 0 & 3
\end{bmatrix} \cdot 
\begin{bmatrix}
    4 \\ 5 \\ 6
\end{bmatrix} = 
\begin{bmatrix}
    4 \\ 10 \\ 18
\end{bmatrix}
\end{align}

\lipsum[4] Lorem ipsum dolor sit amet, consectetur adipiscing elit, sed do eiusmod tempor incididunt ut labore et dolore magna aliqua \cite{szczypiorski2015}, \cite{duqu2011}, \cite{shs2015}, \cite{wozniak2018}, \cite{dcp19}. 

\subsection{Critique of Pure Reason}
\kant[1]

\begin{table}[!h] \label{tab:tabela1} \centering
\caption{Przykładowa tabela.}
\begin{tabular} {| c | c | r |} \hline
    Kolumna 1 & Kolumna 2 & Liczba \\ \hline\hline
    cell1 & cell2 & 60 \\ \hline
    cell4 & cell5 & 43 \\ \hline
    cell7 & cell8 & 20,45 \\ \hline
    \multicolumn{2}{|r|}{Suma:} & 123,45 \\ \hline
\end{tabular}
\end{table}

\kant[2]

\begin{longtable}{| c | m{0.58\linewidth} | r | m{0.1\linewidth} |} 
    \caption{Tabela wielostronicowa.} \\ 
    \hline
    Lp & \multicolumn{1}{c|}{Treść} & \multicolumn{1}{c|}{Kwota} & \multicolumn{1}{m{0.1\linewidth}|}{Wariant opłaty} \\ \hline\hline \endfirsthead
    
    \endfoot
    \hline \endlastfoot
    
    1 & Lorem ipsum dolor sit amet, consectetur adipiscing elit, sed do eiusmod tempor incididunt ut labore et dolore magna aliqua. & 111 111,11 zł & \multicolumn{1}{c|}{WAR1} \\ \hline
    2 & Lorem ipsum dolor sit amet, consectetur adipiscing elit, sed do eiusmod tempor incididunt ut labore et dolore magna aliqua. & 22 222,22 zł & \multicolumn{1}{c|}{WAR1} \\ \hline
    3 & Lorem ipsum dolor sit amet, consectetur adipiscing elit, sed do eiusmod tempor incididunt ut labore et dolore magna aliqua. & 33 333,33 zł & \multicolumn{1}{c|}{WAR1} \\ \hline
    4 & Lorem ipsum dolor sit amet, consectetur adipiscing elit, sed do eiusmod tempor incididunt ut labore et dolore magna aliqua. & 444 444,44 zł & \multicolumn{1}{c|}{WAR1} \\ \hline
    5 & Lorem ipsum dolor sit amet, consectetur adipiscing elit, sed do eiusmod tempor incididunt ut labore et dolore magna aliqua. & 55 555,55 zł & \multicolumn{1}{c|}{WAR1} \\ \hline
    6 & Lorem ipsum dolor sit amet, consectetur adipiscing elit, sed do eiusmod tempor incididunt ut labore et dolore magna aliqua. & 66 666,66 zł & \multicolumn{1}{c|}{WAR1} \\ \hline
    7 & Lorem ipsum dolor sit amet, consectetur adipiscing elit, sed do eiusmod tempor incididunt ut labore et dolore magna aliqua. & 777 777,77 zł & \multicolumn{1}{c|}{WAR1} \\ \hline
    8 & Lorem ipsum dolor sit amet, consectetur adipiscing elit, sed do eiusmod tempor incididunt ut labore et dolore magna aliqua. & 8 888,88 zł & \multicolumn{1}{c|}{WAR1} \\ \hline
    9 & Lorem ipsum dolor sit amet, consectetur adipiscing elit, sed do eiusmod tempor incididunt ut labore et dolore magna aliqua. & 999 999,99 zł & \multicolumn{1}{c|}{WAR1} \\ \hline
    10 & Lorem ipsum dolor sit amet, consectetur adipiscing elit, sed do eiusmod tempor incididunt ut labore et dolore magna aliqua. & 111 111,11 zł & \multicolumn{1}{c|}{WAR2} \\ \hline
    11 & Lorem ipsum dolor sit amet, consectetur adipiscing elit, sed do eiusmod tempor incididunt ut labore et dolore magna aliqua. & 22 222,22 zł & \multicolumn{1}{c|}{WAR2} \\ \hline
    12 & Lorem ipsum dolor sit amet, consectetur adipiscing elit, sed do eiusmod tempor incididunt ut labore et dolore magna aliqua. & 33 333,33 zł & \multicolumn{1}{c|}{WAR2} \\ \hline
    13 & Lorem ipsum dolor sit amet, consectetur adipiscing elit, sed do eiusmod tempor incididunt ut labore et dolore magna aliqua. & 444 444,44 zł & \multicolumn{1}{c|}{WAR2} \\ \hline
    14 & Lorem ipsum dolor sit amet, consectetur adipiscing elit, sed do eiusmod tempor incididunt ut labore et dolore magna aliqua. & 55 555,55 zł & \multicolumn{1}{c|}{WAR2} \\ \hline
    15 & Lorem ipsum dolor sit amet, consectetur adipiscing elit, sed do eiusmod tempor incididunt ut labore et dolore magna aliqua. & 66 666,66 zł & \multicolumn{1}{c|}{WAR2} \\ \hline
    & \multicolumn{1}{r|}{\textbf{Suma:}} & \textbf{7 777 777,77 zł} & 
    \label{table:koszty}
\end{longtable}
\kant[4]

\subsection{Categorical Imperative}
As any dedicated reader can clearly see, the Ideal of practical reason is a representation of, as far as I know, the things in themselves; as I have shown elsewhere, the phenomena should only be used as a canon for our understanding:
% Parametr label ustawia symbol, a leftmargin - wielkość wcięcia.
% Domyślny układ to [---] bez wcięcia, bo tak pan Marcin Woliński powiedział;
% ale ja nie polecam. // AB
\begin{itemize}
    \item Item 1:
    \begin{itemize}[label=---]
        \item item 1.1;
        \item item 1.2;
        \item item 1.3;
    \end{itemize}
    \item Item 2;
    \item Item 3;
    \item Item 4.
\end{itemize}
\kant[2]
\begin{enumerate}
    \item Item 1:
    \begin{enumerate}
        \item item 1.1;
        \item item 1.2:
        \begin{enumerate}
            \item item 1.2.1;
            \item item 1.2.2;
        \end{enumerate}
        \item item 1.3;
    \end{enumerate}
    \item Item 2;
    \item Item 3;
    \item Item 4.
\end{enumerate}

\kant[9]

\subsection{G\"odel's ontological proof}
\kant[9] Lorem ipsum dolor sit amet, consectetur adipiscing elit, sed do eiusmod tempor incididunt ut labore et dolore magna aliqua \cite{benzmuller2014}, \cite{goedel95}, \cite{wang97}, \cite{koons2005}. 
\begin{assumption} \label{ass:1}
    $ [\![ \ \phi \ ]\!] \Longrightarrow [\![ \ P(\phi); \neg P(\phi) \ ]\!]$
\end{assumption}
\begin{axiom}[Dualność] \label{axiom:1}
    $\neg P(\phi) \Leftrightarrow P(\neg \phi)$, równoważnie $P(\phi) \Leftrightarrow \neg P(\neg \phi)$
\end{axiom}
\begin{axiom}[Całkowitość] \label{axiom:2}
    $ \left( P(\phi) \wedge \forall x: \phi(x) \Rightarrow \psi(x) \right) \Rightarrow P(\psi) $
\end{axiom}
\begin{axiom}[Absolutność] \label{axiom:3}
    $ P(\phi) \Rightarrow \Box P(\phi) $
\end{axiom}
\begin{definition} \label{def:1}
    $ G(x) \Leftrightarrow \forall \phi: \left( P(\phi) \Rightarrow \phi(x) \right) $
\end{definition}
\begin{definition} \label{def:2}
    $ \phi \ ess \ x \Leftrightarrow \phi(x) \wedge \forall \psi \left( \psi(x) \Rightarrow \Box \forall y \left( \phi(y) \Rightarrow \psi(y) \right) \right)  $
\end{definition}
\begin{axiom} \label{axiom:4}
    P(G)
\end{axiom}
\begin{lemma} \label{lemma:1}
    $ P(\phi) \Rightarrow \Diamond \exists x : \phi(x) $
\end{lemma}
\begin{proof}
    Dowód pomijamy, bo jest trywialny :)
\end{proof}
\begin{lemma} \label{lemma:2}
    $ \Diamond \exists x : G(x) $
\end{lemma}
\begin{proof}
    Natychmiastowy wniosek z aksjomatu \ref{axiom:4} i lematu \ref{lemma:1}.
\end{proof}
\begin{lemma} \label{lemma:3}
    $ G(x) \Rightarrow G \ ess \ x $
\end{lemma}
\begin{proof}
    Poprzez podstawienie do definicji \ref{def:2}. 
\end{proof}
\begin{definition} \label{def:3}
    $ E(x) \Leftrightarrow \forall \phi \left( \phi \ ess \ x \Rightarrow \Box\ \exists x: \phi(x) \right) $
\end{definition}
\begin{axiom} \label{axiom:5}
    P(E)
\end{axiom}
\begin{theorem}
    $ \Box\ \exists x : G(x) $
\end{theorem}
\begin{proof}
    Na podstawie definicji \ref{def:1}, lematu \ref{lemma:3} i aksjomatu \ref{axiom:5}. 
\end{proof}    % wygodnie jest trzymać
% \newpage % Zaleca się otwieranie rozdziału od nowej strony.
\section{Code listings}
\lipsum[10]
\begin{lstlisting}[language=HTML]
<html>
  <head>
    <title>Hello world!</title>
  </head>
  <body>
    Hello
  </body>
</html>
\end{lstlisting}
\lipsum[11]
\begin{lstlisting}[language=C]
#include <stdio.h>
int main() {
  // printf() displays the string inside quotation
  printf("Hello world!");
  return 0;
}
\end{lstlisting}
\lipsum[12]
 % każdy rozdział w osobnym pliku. 

%--------------------------------------------
% Literatura
%--------------------------------------------
\newpage
\printbibliography

%--------------------------------------------
% Spisy (opcjonalne)
%--------------------------------------------
\newpage

% Wykaz symboli i skrótów.
% Pamiętaj, żeby posortować symbole alfabetycznie
% we własnym zakresie. Ponieważ mało kto używa takiego wykazu, 
% uznałem, że robienie automatycznie sortowanej listy
% na poziomie LaTeXa to za duży overkill. 
% Jest tylko proste i oczywiste makro \acronym, 
% które dodaje postawowe formatowanie. 
% //AB
\vspace{0.8cm}
\section*{Wykaz symboli i skrótów}
\acronym{EiTI}{Wydział Elektroniki i Technik Informacyjnych}
\acronym{PW}{Politechnika Warszawska}

\listoffigures              % Spis obrazków. 
\vspace{1cm}                % vertical space
\listoftables               % Spis tabel. 
\vspace{1cm}                % vertical space
\listofappendices           % Spis załączników

% Załączniki
\appendix

\newpage
\newappendix{Nazwa załącznika 1}
\lipsum[1]

\newpage
\newappendix{Nazwa załącznika 2}
\lipsum[1]

% Używając powyższych spisów jako szablonu,
% możesz tu dodać swój własny wykaz bądź listę, 
% np. spis algorytmów. 

\end{document} % Dobranoc. 

